\chapter{Benign Prostatic Hyperplasia}
\index{Benign Prostatic Hyperplasia}

\section*{Definition and Overview}
Benign prostatic hyperplasia (BPH) is a non-cancerous enlargement of the prostate gland that occurs as men age. The prostate surrounds the urethra, the tube that carries urine from the bladder out of the body. As the gland grows, it can compress the urethra and obstruct the flow of urine, giving rise to a cluster of bothersome urinary symptoms collectively known as lower urinary tract symptoms (LUTS).

BPH is not cancer, does not by itself increase the risk of prostate cancer, and is not life-threatening. However, it can significantly reduce quality of life through sleep disturbance, discomfort, and urinary complications. A range of effective treatments, from lifestyle measures and medication to surgery, means that most men with BPH can achieve good control of their symptoms.

\section*{Epidemiology}
BPH is one of the most common conditions affecting older men, and its frequency rises steeply with age.

\begin{itemize}
    \item \textbf{Histological prevalence:} microscopic enlargement of the prostate is present in around 50 per cent of men in their 50s and up to 90 per cent of men in their 80s
    \item \textbf{Symptomatic disease:} roughly 1 in 4 men develops symptoms that require treatment at some stage
    \item \textbf{Age at onset:} symptoms usually first appear after the age of 40 and become more common and more troublesome with each decade
    \item \textbf{Family history:} men with a first-degree relative affected by BPH have a higher risk
    \item \textbf{Metabolic factors:} obesity, diabetes, and physical inactivity are associated with a greater likelihood of developing symptoms
\end{itemize}

\section*{Aetiology and Pathophysiology}
The precise cause of BPH is not fully understood, but ageing and the male hormone testosterone are both required for the prostate to enlarge.

As men age, testosterone within the prostate is converted by the enzyme 5-alpha reductase into the more potent androgen dihydrotestosterone (DHT). DHT stimulates the cells of the prostate to multiply, so the gland gradually increases in size. The growth mainly affects the transition zone of the gland, the tissue lying immediately around the urethra, so even modest enlargement can narrow the channel and obstruct the outflow of urine.

In response to the increased resistance, the muscle of the bladder wall (detrusor) thickens and must work harder to empty the bladder. Over time the bladder may become less able to generate a strong stream, may develop instability that causes urgency and frequency, and may fail to empty completely. The overall symptom burden reflects a combination of the mechanical obstruction caused by the enlarged gland and the functional changes in bladder muscle, which is why symptom severity does not always match the size of the prostate.

\section*{Clinical Features}
Symptoms of BPH usually develop gradually over years and are grouped under the heading of lower urinary tract symptoms.

\begin{itemize}
    \item \textbf{Weak or interrupted stream:} a diminished force of urine flow, or a flow that starts and stops
    \item \textbf{Hesitancy:} difficulty starting to urinate and the need to strain to pass urine
    \item \textbf{Terminal dribble:} leakage of urine after voiding has apparently finished
    \item \textbf{Incomplete emptying:} a persistent feeling that the bladder is not fully empty
    \item \textbf{Urgency:} a sudden, strong, and difficult-to-control need to urinate
    \item \textbf{Frequency and nocturia:} the need to pass urine more often during the day and to wake at night to urinate
    \item \textbf{Acute urinary retention:} in severe cases, complete inability to pass urine despite a full bladder, which is a medical emergency
    \item \textbf{Complications:} recurrent urinary tract infections, blood in the urine, or the formation of bladder stones in some men
\end{itemize}

\section*{Differential Diagnosis}
Several conditions can produce urinary symptoms very similar to those of BPH.

\begin{itemize}
    \item \textbf{Prostate cancer:} can cause similar symptoms, and the two conditions frequently coexist in older men
    \item \textbf{Urinary tract infection or prostatitis:} inflammation of the bladder or prostate, often with pain, fever, or burning
    \item \textbf{Urethral stricture:} narrowing of the urethra, commonly a consequence of past injury, surgery, or infection
    \item \textbf{Overactive bladder:} urinary urgency and frequency without obstruction of the bladder outlet
    \item \textbf{Neurological conditions:} such as spinal cord injury, multiple sclerosis, or peripheral neuropathy, which impair bladder control
    \item \textbf{Medication side effects:} some cold, allergy, antihistamine, and antidepressant medicines can worsen urinary symptoms or cause retention
    \item \textbf{Diabetes:} can cause increased urine output and bladder nerve dysfunction
\end{itemize}

\section*{When to Seek Medical Care}
See a doctor if urinary symptoms are new, persistent, bothersome, or interfering with sleep and daily activities, or if there is blood in the urine, pain on urination, or a fever. Men with a family history of prostate cancer or who are of an age and risk profile that warrant screening should discuss the role of PSA testing.

\textbf{Contact your local emergency services for:}
\begin{itemize}
    \item Sudden inability to pass urine despite a full bladder, especially with severe lower abdominal pain (acute urinary retention)
    \item High fever with back or loin pain, shaking, vomiting, or confusion, which may indicate a severe kidney infection
    \item Inability to drink, or persistent vomiting
    \item Sudden weakness, chest pain, or difficulty breathing
\end{itemize}

\section*{Diagnosis and Investigations}
\begin{itemize}
    \item \textbf{History:} a detailed discussion of urinary symptoms, often using a validated symptom score such as the International Prostate Symptom Score (IPSS), together with review of medicines, fluid intake, and other health conditions
    \item \textbf{Digital rectal examination (DRE):} the doctor feels the prostate through the rectum to assess its size, shape, and texture; a hard, irregular, or fixed gland raises the possibility of cancer
    \item \textbf{PSA blood test:} prostate-specific antigen levels help estimate the chance of prostate cancer; elevated or rapidly rising levels require further investigation, although PSA can also be raised by BPH, infection, and recent instrumentation
    \item \textbf{Urine tests:} dipstick testing and culture to exclude infection, blood, or other abnormalities
    \item \textbf{Post-void residual volume:} ultrasound measurement of urine left in the bladder after voiding, which indicates the efficiency of bladder emptying
    \item \textbf{Uroflowmetry:} a measurement of the strength and pattern of the urine stream
    \item \textbf{Kidney function tests:} blood tests to check for renal damage from long-standing obstruction
    \item \textbf{Urology referral:} specialist review, cystoscopy, or prostate biopsy if cancer is suspected, investigations are unclear, or complications are developing
\end{itemize}

\section*{Management}
Treatment is guided by how troublesome the symptoms are, the size of the prostate, and whether complications are developing. Options range from simple monitoring through medication to surgery.

\begin{itemize}
    \item \textbf{Watchful waiting:} appropriate for mild symptoms, with regular review, symptom monitoring, and a PSA check
    \item \textbf{Lifestyle measures:} reducing evening fluids, limiting caffeine and alcohol, urinating at set times, and double voiding (urinating, waiting, then urinating again)
    \item \textbf{Alpha-blockers:} such as tamsulosin, relax the smooth muscle of the prostate and bladder neck, improving flow and reducing symptoms; they act quickly but do not shrink the gland
    \item \textbf{5-alpha reductase inhibitors:} such as finasteride and dutasteride, shrink the prostate by blocking the production of DHT; they take several months to work and are most effective for larger glands
    \item \textbf{Combination therapy:} an alpha-blocker together with a 5-alpha reductase inhibitor for men with moderate to severe symptoms and an enlarged prostate
    \item \textbf{Surgery:} for moderate to severe symptoms, failed medication, or complications; transurethral resection of the prostate (TURP) and laser procedures such as HoLEP are performed through the urethra without an external incision
    \item \textbf{Minimally invasive therapies:} including water vapour (Rezum) and prostate artery embolisation, which are options in selected cases
    \item \textbf{Follow-up:} regular monitoring of symptoms, urine flow, PSA, and kidney function
\end{itemize}

\section*{Complications}
\begin{itemize}
    \item \textbf{Acute urinary retention:} sudden inability to urinate, requiring emergency catheterisation
    \item \textbf{Chronic retention:} longstanding incomplete emptying with progressive bladder and kidney damage, sometimes without obvious symptoms
    \item \textbf{Recurrent urinary tract infections} from urine that stagnates in the bladder
    \item \textbf{Bladder stones:} formed from crystals in urine left behind after voiding
    \item \textbf{Kidney damage:} renal impairment from long-standing back-pressure on the upper urinary tract
    \item \textbf{Complications of treatment:} surgical bleeding, infection, narrowing of the urethra, urinary incontinence (uncommon), and retrograde ejaculation, in which semen passes into the bladder rather than out through the penis
\end{itemize}

\section*{Prognosis and Outcomes}
BPH is not a life-threatening condition, but untreated it can significantly impair quality of life and, in severe cases, damage the kidneys. Symptoms may remain stable for years or progress slowly, and progression is not predictable in an individual man.

Many men improve substantially with medication, and surgery relieves obstruction in the great majority of those who need it, often for many years. Kidney damage is uncommon when the condition is monitored and treated appropriately. BPH does not cause prostate cancer, but because both are common in older men, each requires separate assessment and appropriate surveillance.

\section*{Prevention and Self-Care}
BPH cannot always be prevented, but several measures may reduce symptoms and limit their impact.

\begin{itemize}
    \item Maintain a healthy weight and remain physically active
    \item Limit caffeine and alcohol, which can irritate the bladder, particularly in the evening
    \item Avoid drinking large amounts of fluid in the hours before bed
    \item Do not ignore the urge to urinate, and do not delay urination for long periods
    \item Ask a doctor or pharmacist whether any medicines, including decongestants and some antihistamines, could be worsening symptoms
    \item Keep warm, because cold weather can increase urinary symptoms
    \item Attend recommended check-ups so that symptoms, PSA, and kidney function are monitored over time
\end{itemize}

\section*{Sources and Further Reading}
The following organisations provide reliable information on benign prostatic hyperplasia for patients, families, and health professionals.

\begin{itemize}
    \item NHS. \textit{Information on an enlarged prostate and urinary symptoms.} https://www.nhs.uk/conditions/
    \item Mayo Clinic. \textit{Overview of benign prostatic hyperplasia and its treatment.} https://www.mayoclinic.org/
    \item MSD Manual. \textit{Professional information on benign prostatic hyperplasia.} https://www.msdmanuals.com/
    \item NICE. \textit{Clinical guidance on the assessment and management of lower urinary tract symptoms in men.} https://www.nice.org.uk/
    \item MedlinePlus. \textit{Health information on benign prostatic hyperplasia for patients.} https://medlineplus.gov/
\end{itemize}
